% =====================================================================
%  Seksjon: Nye co-simuleringsscenarioer på den reduserte modellen mot
%           hele LEOGO-nettet (driftshendelse, frekvensstøtte-avveining,
%           resonant prosesslast, resonanskurve).
%  Driver:  casestudies/dyn_sim/test_WT_LEOGO_tower_sim.py
%  Analyse: casestudies/dyn_sim/_analyze_gt_trip.py
%           casestudies/dyn_sim/_analyze_freq_support.py
%           casestudies/dyn_sim/_analyze_resonance.py
%           casestudies/dyn_sim/sweep_ss_resonance.py
%  Data:    results/tower_test/{gt_trip,gt_trip_support,resonance_0234,
%                               ss_resonance_curve}.{csv,png}
%  Preamble: amsmath, amssymb, graphicx, booktabs, siunitx
%           + \DeclareSIUnit{\pu}{pu}   (per-unit, used by \SI{}{\pu})
% =====================================================================

\section{New Co-Simulation Scenarios on the LEOGO Network}
\label{sec:new-tests}

With the tower modes calibrated (\S\ref{sec:simplified-tower-model}) and the
shaft-torque drive path established (\S\ref{sec:shaft-excitation}), the reduced
\texttt{WindTurbineTower} is exercised against the full LEOGO microgrid in a
common test harness (\texttt{test\_WT\_LEOGO\_tower\_sim.py}). The turbine
replaces the original LEOGO WTG converter at \texttt{Busbar WTG1 LV} through a
UIC interface, and the three gas-turbine generators, their governors and AVRs,
and the platform loads are simulated as usual with a fixed communication step of
\SI{0.01}{\second}. The harness applies disturbances at the main gas-turbine bus
--- a smooth (ramped) load step, a sustained sinusoidal load modulation, or a
short circuit --- and logs the grid frequency, the turbine electrical variables,
and both tower accelerations. Four scenarios are reported, each targeting a
different aspect of the electro-mechanical interaction.

\subsection{A discrete grid event: the tower as a mechanical band-pass filter}
\label{subsec:nt-gttrip}

The first scenario is a realistic single disturbance: a step change in platform
load, representative of a large consumer switching or a generator trip, applied
as a smooth shunt load at the main bus. The step is broadband --- most of its
spectral energy lies near DC --- yet the tower side-to-side response is a clean
decaying oscillation at \(f_{\mathrm{ss}}\approx\SI{0.234}{\hertz}\), the modal
frequency, not at the (much slower) frequency of the grid transient itself. The
step acts as an impulsive ``kick'' that excites the free side-to-side ring-down
of \eqref{eq:stm-tau}, which decays with \(\tau_{\mathrm{ss}}\approx
\SI{13.6}{\second}\). The fore-aft mode barely responds: the ratio of peak
side-to-side to peak fore-aft acceleration is of order \(\sim\!40\), confirming
that even a broadband electrical event is filtered by the tower into an almost
purely side-to-side motion. This is the mechanical-band-pass picture of
\S\ref{subsec:se-selectivity}, seen for a discrete event
(Figure~\ref{fig:nt-gttrip}).

\begin{figure}[t]
  \centering
  % \includegraphics[width=0.8\linewidth]{fig/gt_trip.png}
  \caption{Discrete grid event (reduced model): a broadband platform-load step
  produces a decaying side-to-side ring-down at the tower modal frequency
  (\SI{0.234}{\hertz}), while the fore-aft mode stays essentially flat. The tower
  band-pass-filters the electrical disturbance onto its side-to-side mode.}
  \label{fig:nt-gttrip}
\end{figure}

\subsection{The frequency-support trade-off}
\label{subsec:nt-tradeoff}

The headline scenario quantifies the central tension of the thesis: the wind
turbine's frequency-support droop (\S\ref{sec:outer-droop}) helps the grid but
loads the tower. A \emph{temporary} \SI{5}{\mega\watt} load step is applied at
the main bus (smooth \SI{2}{\second} ramps, held for \SI{50}{\second}, then
removed), and the co-simulation is run twice with the identical de-loaded
operating point: once with the droop off and once with it on
(\(K_{\mathrm{droop}}=\SI{0.75}{\pu\per\hertz}\), \SI{10}{\percent} headroom).
Because the tower modes are forward-coupled (\S\ref{subsec:stm-forward}), the two
runs differ only in the droop action, so the comparison is clean.

The results are collected in Table~\ref{tab:nt-tradeoff}. On the grid side the
droop helps: it cuts the settled (plateau) frequency dip by \SI{20}{\percent}
(from \SI{0.089}{\hertz} to \SI{0.071}{\hertz}) and the transient dip by
\SI{15}{\percent}, by injecting an extra \(\approx\SI{0.008}{\pu}\) of active power
from the headroom reserve. On the mechanical side the same injection is the
penalty: the side-to-side acceleration grows by a factor \(1.72\) in both peak
and RMS. The fore-aft response stays negligible. The interaction is thus not
merely academic --- providing a genuine grid service (frequency support) directly
and measurably increases the resonance-prone tower loading, a coupling that a
turbine-level or grid-level design would need to weigh against each other.

\begin{table}[t]
  \centering
  \caption{Frequency-support trade-off (reduced model): temporary
  \SI{5}{\mega\watt} load step, droop off versus on
  (\(K_{\mathrm{droop}}=\SI{0.75}{\pu\per\hertz}\), \SI{10}{\percent} headroom).
  The droop reduces the grid-frequency dip but increases the tower side-to-side
  loading.}
  \label{tab:nt-tradeoff}
  \begin{tabular}{lccc}
    \toprule
    Quantity & Off & On & Ratio on/off \\
    \midrule
    Transient freq.\ dip [\si{\hertz}]     & 0.0897 & 0.0760 & 0.85 \\
    Plateau freq.\ dip [\si{\hertz}]       & 0.0893 & 0.0714 & 0.80 \\
    WT power, plateau [\si{\pu}]            & 0.1432 & 0.1513 & --- \\
    \midrule
    SS accel.\ peak [\si{\metre\per\second\squared}] & 0.0137 & 0.0235 & 1.72 \\
    SS accel.\ RMS  [\si{\metre\per\second\squared}] & 0.0028 & 0.0049 & 1.72 \\
    FA accel.\ peak [\si{\metre\per\second\squared}] & 0.0005 & 0.0009 & --- \\
    \bottomrule
  \end{tabular}
\end{table}

\begin{figure}[t]
  \centering
  % \includegraphics[width=0.8\linewidth]{fig/gt_trip_support.png}
  \caption{Frequency-support trade-off (reduced model): grid frequency (top), WT
  active power (middle) and tower side-to-side acceleration (bottom) for the
  temporary load step, with the droop off versus on. The droop lifts the power
  and shrinks the frequency dip, but rings the side-to-side mode
  \(\approx\!1.7\times\) harder.}
  \label{fig:nt-tradeoff}
\end{figure}

\subsection{Resonant excitation by a cyclic process load}
\label{subsec:nt-resonance}

Where the load step is broadband and only \emph{rings} the tower down, this
scenario supplies a \emph{sustained} excitation in the tower band and drives the
side-to-side mode into genuine resonance. It represents a realistic islanded-grid
situation: a large cyclic mechanical load --- a reciprocating compressor or a
beam / sucker-rod pump running near 14 strokes per minute
(\(=\SI{0.234}{\hertz}\)) --- appears as a periodic active-power oscillation at
the main bus. A \(\pm\SI{5}{\mega\watt}\) sinusoidal load at \SI{0.234}{\hertz}
is switched on at \(t=\SI{20}{\second}\), producing a grid-frequency oscillation
of \(\approx\SI{67}{\milli\hertz}\) amplitude at that frequency.

Because \SI{0.234}{\hertz} coincides with the side-to-side natural frequency, the
mode is driven resonantly through the open path \eqref{eq:se-path} and its
amplitude \emph{builds up} over roughly \(3\text{--}4\,\tau_{\mathrm{ss}}\approx
\SI{45}{\second}\) to a large sustained value, rather than decaying as in
\S\ref{subsec:nt-gttrip} (Figure~\ref{fig:nt-resonance}). The settled side-to-side
lock-in acceleration reaches
\(A_{\mathrm{ss}} = \SI{0.191}{\metre\per\second\squared}\) --- about four times
the peak of the broadband step --- while the fore-aft lock-in is only
\(A_{\mathrm{fa}} = \SI{0.0063}{\metre\per\second\squared}\), a side-to-side/
fore-aft ratio of \(\approx\!30\). The build-up versus ring-down contrast is the
defining signature of resonant, as opposed to broadband, excitation.

\begin{figure}[t]
  \centering
  % \includegraphics[width=0.8\linewidth]{fig/resonance_0234.png}
  \caption{Resonant excitation (reduced model): a \(\pm\SI{5}{\mega\watt}\)
  cyclic process load at \SI{0.234}{\hertz} (top) drives a
  \(\approx\SI{67}{\milli\hertz}\) grid-frequency oscillation (middle), and the
  side-to-side mode builds up over \(\sim\!\SI{40}{\second}\) to a sustained
  amplitude while the fore-aft mode stays flat (bottom, with the side-to-side
  amplitude envelope dashed).}
  \label{fig:nt-resonance}
\end{figure}

\subsection{The resonance curve: frequency selectivity of the coupling}
\label{subsec:nt-curve}

Sweeping the cyclic-load frequency from \SIrange{0.200}{0.300}{\hertz} and
measuring the settled side-to-side lock-in amplitude at each frequency
(\texttt{sweep\_ss\_resonance.py}) traces out the resonance curve of
Table~\ref{tab:nt-curve} and Figure~\ref{fig:nt-curve}. It is a clean, sharp
Lorentzian peak: the side-to-side amplitude climbs from
\SI{0.051}{\metre\per\second\squared} at \SI{0.200}{\hertz} to a maximum of
\SI{0.190}{\metre\per\second\squared} at exactly \SI{0.234}{\hertz} --- matching
both the calibrated modal frequency and the dedicated resonant run
(\S\ref{subsec:nt-resonance}) --- and falls away symmetrically above it. The
half-power width is only \(\approx\SI{0.026}{\hertz}\), a quality factor
\(Q\approx 9\) consistent with \(\zeta_{\mathrm{ss}}=\num{0.05}\)
(\(Q=1/2\zeta=10\)). The grid-frequency drive is essentially flat
(\(\approx\SI{67}{\milli\hertz}\)) across the whole band, so the peak is the
tower's mechanical resonance and not a feature of the excitation. The fore-aft
amplitude stays below \SI{0.0065}{\metre\per\second\squared} everywhere,
including where the load frequency crosses the (nearly coincident) fore-aft
frequency --- once more the drive path, not the frequency, is what selects the
side-to-side mode.

The practical reading is the frequency selectivity itself: a cyclic platform load
produces significant tower loading only if its fundamental or a harmonic falls
within the narrow \(\approx\SI{0.026}{\hertz}\) window around
\SI{0.234}{\hertz}; outside it the coupling is negligible. This both bounds the
risk and identifies precisely which operational frequencies to avoid.

\begin{table}[t]
  \centering
  \caption{Reduced-model resonance curve: settled side-to-side and fore-aft
  lock-in acceleration versus cyclic-load frequency (\(\pm\SI{5}{\mega\watt}\)).
  The bold row marks the side-to-side mode.}
  \label{tab:nt-curve}
  \begin{tabular}{lcc}
    \toprule
    \(f\) [\si{\hertz}] &
      \(A_{\mathrm{ss}}\) [\si{\metre\per\second\squared}] &
      \(A_{\mathrm{fa}}\) [\si{\metre\per\second\squared}] \\
    \midrule
    0.200          & 0.0507 & 0.0009 \\
    0.216          & 0.0948 & 0.0015 \\
    0.226          & 0.1538 & 0.0031 \\
    0.230          & 0.1780 & 0.0052 \\
    \textbf{0.234} & \textbf{0.1903} & \textbf{0.0064} \\
    0.238          & 0.1828 & 0.0059 \\
    0.242          & 0.1623 & 0.0038 \\
    0.248          & 0.1297 & 0.0015 \\
    0.255          & 0.1026 & 0.0011 \\
    0.280          & 0.0578 & 0.0005 \\
    0.300          & 0.0440 & 0.0003 \\
    \bottomrule
  \end{tabular}
\end{table}

\begin{figure}[t]
  \centering
  % \includegraphics[width=0.8\linewidth]{fig/ss_resonance_curve.png}
  \caption{Reduced-model resonance curve: settled side-to-side (red) and fore-aft
  (green) acceleration amplitude versus the frequency of a \(\pm\SI{5}{\mega\watt}\)
  cyclic grid load. The side-to-side response peaks sharply at \SI{0.234}{\hertz}
  (\(Q\approx 9\)); the fore-aft response stays flat, even where the load
  frequency crosses the fore-aft mode.}
  \label{fig:nt-curve}
\end{figure}

\subsection{Summary}
\label{subsec:nt-summary}

Taken together the four scenarios trace the interaction from mechanism to
operational consequence. A discrete grid event shows the tower band-pass-filtering
a broadband disturbance onto its side-to-side mode
(\S\ref{subsec:nt-gttrip}); a sustained cyclic load shows the same mode driven
into resonant build-up (\S\ref{subsec:nt-resonance}); the sweep shows the coupling
is sharply frequency-selective about \SI{0.234}{\hertz} (\S\ref{subsec:nt-curve});
and the frequency-support scenario shows that acting on the grid with the turbine
directly aggravates that tower loading (\S\ref{subsec:nt-tradeoff}). All four run
in real time on the reduced model against the full LEOGO network, complementing
the high-fidelity FMU evidence of \S\ref{sec:tower-ss-resonance} with a fast,
fully closed electrical-to-mechanical demonstration.
